\documentclass[12pt,oneside]{article}

%%%%%%%%%%%%%%%%%%%%%%%%%%%%
%%   Zusaetzliche Pakete  %%
%%%%%%%%%%%%%%%%%%%%%%%%%%%%
\usepackage{acronym}
\usepackage{enumerate}
\usepackage{a4wide}
\usepackage{fancyhdr}
\usepackage{graphicx}
\usepackage{palatino}
\usepackage{blindtext}
\usepackage{multirow}
\usepackage[section]{placeins}
\usepackage[ruled,longend]{algorithm2e}
\usepackage{float}
\usepackage{amsmath}
\usepackage{makecell}
\usepackage{tabularx}
\usepackage{amssymb}
\usepackage{subfiles}


%folgende Zeile auskommentieren für englische Arbeitenhttps://www.overleaf.com/project/62ecf4bc5c4c340d9e1ba297
\usepackage[ngerman]{babel}

\usepackage[bookmarks]{hyperref}
\usepackage[T1]{fontenc}
\usepackage[utf8]{inputenc}
\usepackage[a-1b]{pdfx}
\usepackage[justification=centering]{caption}
%\usepackage[style=unsrt,natbib=true,backend=biber]{biblatex}
\usepackage{csquotes}
\usepackage{url}

%%%%%%%%%%%%%%%%%%%%%%%%%%%%%%
%% Definition der Kopfzeile %%
%%%%%%%%%%%%%%%%%%%%%%%%%%%%%%

\pagestyle{fancy}
\fancyhf{}
\cfoot{\thepage}
\setlength{\headheight}{16pt}

%%%%%%%%%%%%%%%%%%%%%%%%%%%%%%%%%%%%%%%%%%%%%%%%%%%%%
%%  Definition des Deckblattes und der Titelseite  %%
%%%%%%%%%%%%%%%%%%%%%%%%%%%%%%%%%%%%%%%%%%%%%%%%%%%%%

\newcommand{\HSFTitle}[8]{

  \thispagestyle{empty}
\begin{center}
    \includegraphics[width=0.8\textwidth]{logo.eps} \\
    \vspace*{\stretch{1}}
    \end{center}

  %\vspace*{\stretch{1}}
  {\parindent0cm
  \rule{\linewidth}{.7ex}}
  \begin{center}
    \vspace*{\stretch{1}}
    \sffamily\bfseries\Huge
    #1\\
    \vspace*{\stretch{1}}
    \sffamily\bfseries\large
    #3
    \vspace*{\stretch{1}}
  \end{center}
  \rule{\linewidth}{.7ex}

  \vspace*{\stretch{2}}
  \begin{center}
    \Large #2 am #5 der HAW Fulda \\
    \vspace*{\stretch{1}}

    \large Matrikelnummer:  #4 \\[1mm]
    \large Dozent:  #7 \\[1mm]
    \large Fachbetreuer:  #8 \\[1mm]

    \vspace*{\stretch{1}}
    \large Eingereicht am #6
  \end{center}
}


%%%%%%%%%%%%%%%%%%%%%%%%%%%%
%%  Beginn des Dokuments  %%
%%%%%%%%%%%%%%%%%%%%%%%%%%%%



\begin{document}
\makeatletter
  \HSFTitle
      {Projektmanagement: Maßnahmen zur Optimierung der Nachrichtenverteilung in Gruppen und Organisationseinheiten }        % Titel der Arbeit
      {Hausarbeit} % Typ der Arbeit
      {Nicolas Will, Julian Dreuth, Leon Lux}          % Vor- und Nachname des Autors
      {?, ?, 1420067}
      {Fachbereich AI}  % Name des FBs
      {12.11.2024}        % Tag der Abgabe
      {Frank Breitenbach}     % Name des Erstgutachters
      {Daniel Brenzel}    % Name des Zweitgutachters
  \clearpage

\lhead{}
\pagenumbering{Roman}
    \setcounter{page}{1}

%%%%%%%%%%%%%%%%%%%%%%%%%%%%
%%  Kurzzusammenfassung   %%
%%%%%%%%%%%%%%%%%%%%%%%%%%%%
\clearpage
%
\markboth{Abstract}{Abstract}
\section*{Abstract}

Die folgende Arbeit beschäftigt sich mit der Anwendung von im Modul AI1021 erlernten Techniken des Projektmanagements anhand eines Projektes im 
Bereich J3 der Hessischen Zentrale für Datenverarbeitung (HZD). Ziel ist die Analyse und Optimierung der bereichsinternen Kommunikation mit dem Ziel, 
eine neue Messenging-Lösung zu finden, die bisher verwendete Lösungen ersetzen kann. Die schriftliche Ausarbeitung legt dabei ihren Fokus auf 
das Projektmanagement und Projektcontrolling. Dazu werden im Folgenden die durchlaufenen Phasen des Projektes dokumentiert und evaluiert, um als 
abschließenden Schritt gewonnene Erkenntnisse aufzuzeigen sowie ein Fazit zu ziehen.

\markboth{Sperrvermerk}{Sperrvermerk}
\section*{Sperrvermerk}
Die vorliegende Arbeit enthält Interna der HZD und ist daher streng vertraulich zu behandeln. Sie ist ausschließlich als Prüfungsvorlage bestimmt und weder
fremden Dritten zugänglich gemacht, noch ohne Einverständnisse der HZD vervielfältigt werden.


\clearpage
\tableofcontents
\clearpage

\addcontentsline{toc}{section}{\listfigurename}
\listoffigures

\addcontentsline{toc}{section}{\listtablename}
\listoftables
\clearpage


%%%%%%%%%%%%%%%%%%%%%%%%%%%%
%%  Einstellungen  %%
%%%%%%%%%%%%%%%%%%%%%%%%%%%%
\cleardoublepage
\pagenumbering{arabic}
    \setcounter{page}{1}
\lhead{\nouppercase{\leftmark}}

%%%%%%%%%%%%%%%%%%%%%%%%%%%%
%%  Hauptteil  %%
%%%%%%%%%%%%%%%%%%%%%%%%%%%%
\section{Einleitung}

\section{Projektinitiierung}
\subfile{projektinitieerung.tex}

\section{Planungsphase}

\section{Projektdurchführung und Controlling}

\section{Projektabschluss}



%%%%%%%%%%%%%%%%%%%%%%%%%%%%
%% Literaturverzeichnis wird
%% automatisch eingefügt
%%%%%%%%%%%%%%%%%%%%%%%%%%%%
\clearpage
\bibliographystyle{unsrt}
\bibliography{literatur}

%\lhead{}
%\printbibliography
%\addcontentsline{toc}{section}{\bibname}
\appendix


\end{document}
